\section{Evaluation}
\label{sec:eval}
We ask whether consumer windows retain residual obligations where a source
window loses them (RQ1); whether versions and records remain stable through
failure, retry, repair, and uncertain responses (RQ2); whether execution
respects artifact boundaries and CI gates (RQ3); and what planning costs appear
as graph, intent, and history inputs vary (RQ4).
The mechanism and companion campaigns use a pinned released Linux ARM64
binary; the historical campaign uses an earlier binary. Exact versions,
hashes, environments, and protocols remain in their separate campaign records.
None measures the later release-candidate extensions.

\subsection{Study Design and Evidence Boundaries}
\label{sec:eval:diff}
The prospective mechanism campaign crosses five four-package DAGs (chain,
diamond, fan-out, disconnected pairs, and a chain with shortcuts), stable and
beta baselines, four fault targets, and three injection boundaries: before
build completion, before publication, and after an artifact is written but the
test publisher returns failure. The resulting 120 fault episodes each repeat
the fault, repair it, and run once more to test discharge; ten no-fault and ten
foreign-identity controls bring the matrix to 140 episodes. The last boundary
is an explicit adapter response, not a process crash. Build and publish
concurrency are one; seed 20260910 fixes the order.

For RQ1, a finite model projection enumerates ordered four-node DAGs, each
single origin, reachable completion sets, and failure subsets. Against the same
states and observed episode completions, a computed counterfactual removes a
unit when its source completes. This is a mechanism ablation, not a modified
binary or competitor. For RQ2, every episode checks the residual plan against
the original obligations minus completion tags, preserves target versions,
and inspects immutable version/content identity. The local publisher may accept
identical content after an uncertain response and must reject foreign content.

Complementary evidence uses a restricted reference model on 300 real Git
histories with 5--20 packages, layered DAGs, 4--12 intent units, and mixed tag
states; recovery graphs contain 20 or 100 packages. Verdaccio and Docker
registries exercise the mixed chains; CI injects failing and missing tests.
Historical data cover a broader but imperfect oracle profile,
reconstructed policy replay, pinned tool configurations, and eight scale
cells. Every campaign retains its own binary, environment, and unit: model
states, episodes, histories, fixture executions, and verifier predicates are
reported separately. The deterministic matrices do not warrant $p$-values or
a pooled quality score.

\begin{table}[t]
\centering
\caption{Questions, interventions, counting units, and principal outcomes.}
\label{tab:rqs}
\scriptsize
\setlength{\tabcolsep}{2pt}
\begin{tabular}{@{}lp{1.65cm}p{1.45cm}p{3.65cm}@{}}
\toprule
RQ & Intervention & Unit & Outcome \\
\midrule
1 & finite projection; source-window counterfactual & state, transition,
episode & 0 projection violations; counterfactual loses 260 states and
111 obligations in 84/120 fault episodes \\
2 & faults, repair, identity conflict & 140 episodes & 140 pass; 10/10
foreign artifacts refused \\
3 & real registries; CI gates & 60 executions; 3 CI repetitions & configured
barriers and test gates distinguish intended pass/fail paths \\
4 & graph/history factors & 8 process cells & .017~s at 100 layered packages;
19.272~s at 10,000 \\
\bottomrule
\end{tabular}
\end{table}

\subsection{RQ1: Consumer-Window Mechanism}
The finite projection covers 64 ordered four-node DAGs and four single origins:
256 source instances yield 772 reachable completion states, 2,270
failure-subset transitions, and 1,054 progress transitions. It finds no
violation of exact residual obligations, no widening, no double ledger
delivery, or strict decrease when progress occurs. The source-window mutant
loses obligations in 260 of those states. This is an executable stable-channel,
one-intent projection, not a proof of the full model.

On binary-derived completion states, the same counterfactual loses 111
obligations in 84 of the 120 fault episodes: 28/40 episodes at each of the
build, prepublication, and post-write boundaries, balanced equally between
stable and beta. The ten no-fault controls lose none. All ten foreign-identity
controls also expose one loss each, reported separately because they test a
different intervention. \emph{Answer to RQ1.} Consumer-relative windows retain
the expected residual set in the explored states; source completion is
insufficient evidence that every consumer has discharged the intent.

\subsection{RQ2: Retry and Identity Stability}
All 140 mechanism episodes passed their predeclared outcomes. In the 130 fault
or foreign-identity episodes, repeating the adverse condition preserved the
plan, target versions, and completed tags; repair produced an empty plan. All
140 final invocations were no-ops with stable tags, and the adapter recorded
one artifact write per obligation. Each of the 40 cases at each fault boundary
passed, as did all 70 stable and 70 beta cases. The ten foreign artifacts were
refused without a new completion tag, after which restoring the intended
content allowed recovery. This supports the supplied immutable local adapter,
not arbitrary registry publishers or exactly-once shell effects.

\emph{Answer to RQ2.} Across the designed graph, channel, target, and fault
factors, recorded progress survived repeated failure without version drift;
identity reconciliation distinguished prior intended writes from conflicting
content.

\paragraph{Reference agreement and specification regressions.}
\label{sec:eval:fresh}
All 300 fresh core instances produced accepted matching plans, with no
differences, refusals, or invalid execution. Of 1,500 historical full-profile
histories, 997 produced accepted plans and 503 produced complete E195
refusals for proposed regressive versions. Among accepted plans, 546
comparisons were clean; four refused dossiers also had no row difference,
giving 550 matching diagnostic comparisons. Across the profile, 1,914 package
rows differed. Rejected diagnostic output is not an executable release plan.

The 1,914 differing rows are retained for diagnosis. They show that the
restricted oracle cannot support a full refinement claim, especially for
prerelease and channel behavior. Six historical specification vectors add
selected discharge, suppression, channel, and source-exclusion cases without
covering the complete language.

\paragraph{Earlier recovery and historical replay.}
All 20 companion episodes made progress on every run and converged within
four runs; the verifier inspected 251 command records. Historical evidence
contains two isolated orphan repetitions and 24 broader episodes, which also
converged within four runs with zero structural G3/G5 check failures. These
checks concern tested schedules and predicates. They do
not establish the liveness premise under permanent registry failure, nor do
they make an unrecorded external write unambiguous.

Historical fault cells retain orphan and competing-push cases. The latter
distinguish completed release tags from a concurrent source change; they test
selected interleavings rather than a systematic concurrency sweep.

Five-release historical replay matched 36 of 42 versions in
\code{google-cloud-node} and 7 of 64 in \code{conventional-changelog}. Those
projects did not use the modeled intent language. Their layouts and release
policies were reconstructed for the replay, and mismatches may reflect that
reconstruction as well as different release conventions. Consequently the
experiment demonstrates limits of policy transfer; it is not evidence that
our protocol reproduces arbitrary tools' historical version decisions.

\subsection{Historical Comparisons and Negative Evidence}
\label{sec:compare}
The historical refusal and competing-push fixtures pin Lerna, Nx, and
Changesets releases and configurations in their provenance. Lerna and Nx tagged the
denied version before its failed upload, while Changesets left it untagged
and recovered by querying npm. The observed propagation policies were
unconditional for Lerna/Nx, range-gated for Changesets, and per-commit for our
protocol. These are observations of the pinned configurations, not universal
properties of every mode those tools expose.
In competing-push fixtures, compared runs left tags outside branch ancestry;
our historical run merged before its final push and planned the new commit.
Release-please's tag timing comes from documentation, not an executed
comparison~\cite{releaseplease}. No cross-tool speed or aggregate superiority
claim follows.

\input{mixed-results}

\subsection{RQ4: Planning Cost}
Eight historical cells report seconds and peak MiB. For layered graphs with
200 units, medians rise from .017/15 at 100 packages (185 edges), through
.250/21 at 1,000 (2,002 edges), to 19.272/92 at 10,000 (19,919 edges). A
10,000-package chain records 19.354/118; 1,000-package chain and star cells
record .254/26 and .244/19. At 1,000 layered packages, 1,000 units give
.290/26 and 5,000 padding commits give .664/30. These are contended-host
observations, not confidence intervals or a cross-tool benchmark.
\emph{Answer to RQ4.} Observed costs increase sharply between the 1,000- and
10,000-package cells; the design does not isolate a general asymptotic or
end-to-end bound from history and output costs.
